\BeleskeID{09_1-s040}
% element: 09_1-s040:n01
Ako se $V_{IL}$ nalazi u ovoj oblasti, nagib karakteristike zamenjuje se vrednošću $-1$. Tako nastaje druga relacija, koja ulaz izražava preko izlaza. Njena zamena u prvu jednačinu daje
\[k_{n,L}(V_{DD}-V_O)(V_{DD}+2\Delta-V_O)=k_{n,D}\left(V_{Tn,D}+\frac{k_{n,L}}{k_{n,D}}(V_{DD}+\Delta-V_O)-V_{Tn,D}\right)^2\]
Članovi $V_{Tn,D}$ u zagradi se poništavaju:
\[k_{n,L}(V_{DD}-V_O)(V_{DD}+2\Delta-V_O)=\frac{k_{n,L}^2}{k_{n,D}}(V_{DD}+\Delta-V_O)^2\]
Množenjem sa $k_{n,D}/k_{n,L}$ dobija se poslednji oblik na slajdu. Od tog trenutka jednačina ima samo nepoznati izlazni napon u tački praga.
